
Results provide evidence that both components of bidirectional alignment---user learning and AI responsiveness---exist in RLHF-trained conversational systems, though with weaker effects than initially hypothesized.

\subsubsection{Interpreting the Evidence}

\textbf{User Learning: Small Effect, Real Signal}

The negative reformulation slope (mean = -0.0214, p=0.012, d=-0.246) confirms users adapt their query strategies during conversations. This learning manifests as decreasing reformulation rates. The small effect size indicates learning is subtle rather than dramatic, but the statistical significance across 88 conversations suggests the signal is real.

The effect is small likely because reformulation is a coarse-grained behavioral measure. Users may learn AI capabilities without exhibiting reformulation patterns---for instance, by asking entirely different questions based on gained understanding. The metric captures only one manifestation of learning.

The median slope = 0 finding reveals learning is not universal. Most conversations show flat slopes, while a minority show strong negative slopes. This heterogeneity could reflect individual differences in learning aptitude, variation in task difficulty, or engagement decay confounding true learning. Future work must distinguish between learning (negative slope predicts success) and disengagement (negative slope reflects giving up) by stratifying on conversation outcomes.

\textbf{AI Responsiveness: Below Threshold, Above Zero}

The diversity correlation (r=0.396, p<0.001) demonstrates the RLHF-trained AI system exhibits behavioral responsiveness---response vocabulary varies with query vocabulary. However, the correlation falls below the preregistered threshold (r>0.4).

This threshold failure is attributed to measurement limitations rather than absence of responsiveness. Distinct-1 captures only lexical diversity, not semantic diversity. An AI system that responds to diverse queries with semantically varied responses but overlapping vocabulary would show low Distinct-1 correlation despite genuine responsiveness. Recomputing with semantic diversity metrics would test whether semantic responsiveness is stronger.

The temporal dynamics finding---correlation strengthens from r=0.04 (2-3 turns) to r=0.19 (6+ turns)---provides evidence for co-adaptation. If responsiveness were a static property of the RLHF policy, correlation should be constant across conversation lengths. The observed strengthening suggests responsiveness emerges through extended interaction. However, this could also reflect a confound (longer conversations accumulate more vocabulary variation), requiring partial correlation analysis to isolate true temporal effects.

\textbf{Bidirectional Alignment Framework: Partially Validated}

Results validate the theoretical foundation of bidirectional alignment but leave the central hypothesis---coupling predicts conversation quality beyond individual metrics---untested. Both components exist (learning and responsiveness), but whether they interact to produce better conversations is unknown.

The planned coupling analysis will test whether correlation between reformulation slope and diversity correlation within conversations predicts helpfulness ratings better than slope alone or diversity alone. If coupling adds predictive power, it validates bidirectional alignment as a meaningful construct. If not, learning and responsiveness may be independent contributors to conversation quality rather than synergistic co-adaptation.

\subsubsection{Limitations}

\textbf{L1: Small Sample for Reformulation Analysis (n=88)}

Reformulation slope analysis uses only 88 conversations, representing 4.4\% of the sampled dataset. This small sample results from conservative reformulation detection thresholds prioritizing precision over recall. While the sample is sufficient for statistical significance (p=0.012), it may not generalize to the full HH-RLHF corpus or other datasets.

Mitigation: Expanding to the full 161k conversation dataset and relaxing detection thresholds would test effect robustness. Cross-dataset validation would assess generalization beyond HH-RLHF.

\textbf{L2: Threshold Failure for Responsiveness}

The r=0.396 < 0.4 threshold failure weakens the claim that AI exhibits strong responsiveness. While the correlation is statistically significant and consistent with medium effect sizes in behavioral research, it falls short of the preregistered threshold.

Mitigation: (1) Recompute with semantic diversity metrics to test whether lexical diversity underestimates effect size. (2) Apply the planned helpfulness filter (conversations with ratings > median) to test whether responsiveness is stronger in high-quality conversations. (3) Revise the hypothesis to reflect weak-to-medium responsiveness ($r\approx 0.35-0.40)$.

\textbf{L3: Reformulation Detection Not Validated}

The reformulation detection heuristic (SBERT similarity > 0.7 AND edit distance > 0.3) has not been validated against human annotations. False positives would artificially inflate slope estimates, while false negatives would underestimate slopes. The net effect is unknown.

Mitigation: Human annotation study on a 100-conversation subset would quantify precision and recall. High precision/recall (>0.8) would validate the heuristic; low precision/recall would require recomputing slopes with a validated detector.

\textbf{L4: Correlation vs Causation}

Results show correlation (reformulation decreases over turns, response diversity correlates with query diversity) but do not establish causation. Correlation could reflect confounds or reverse causation.

Mitigation: Causal intervention study where AI responsiveness is experimentally manipulated would establish causal direction. Alternatively, instrumental variable analysis using conversation metadata could estimate causal effects from observational data.

\textbf{L5: Coupling Hypothesis Untested}

The core hypothesis---bidirectional alignment coupling predicts helpfulness---remains untested. Results establish that components exist but not that they interact productively.

Mitigation: Complete the planned coupling analysis: regress helpfulness ratings on coupling strength while controlling for slope alone and diversity alone. Significant coupling coefficient would validate the interaction hypothesis.

\subsubsection{Broader Impact}

\textbf{Implications for Alignment Measurement}

This work suggests alignment should be evaluated not just through static policy quality but through conversation dynamics. Two AI systems with identical per-turn response quality may differ in bidirectional alignment---one elicits user learning and exhibits responsiveness, while the other does not. Standard RLHF evaluation would rate them equally; bidirectional metrics would differentiate them.

Monitoring reformulation slopes and diversity correlations in production logs could detect alignment degradation that aggregate helpfulness ratings miss. If users increasingly reformulate queries without learning (flat slopes despite high reformulation rates), this signals the AI is not adapting effectively despite potentially high per-turn quality.

\textbf{Implications for Training Methods}

Current RLHF optimizes AI policy to maximize human preference predictions, treating users as static evaluators. Results suggest users are not static---they learn during conversations. Future alignment training could incorporate bidirectional objectives:

\begin{itemize}
\item User learning objective: Encourage responses that reduce future reformulation rates.
\item AI responsiveness objective: Train policies that adjust output characteristics based on query patterns.
\item Coupling objective: Maximize correlation between user learning rate and AI responsiveness, directly optimizing for co-adaptation.
\end{itemize}

These objectives would shift RLHF from unidirectional preference matching to bidirectional co-adaptation.

\textbf{Societal Considerations}

Optimizing for bidirectional alignment raises concerns about user dependency. If AI systems are trained to be maximally responsive to user patterns, users may become reliant on AI communication styles and struggle with less adaptive systems.

Balancing bidirectional alignment optimization with user autonomy is recommended. AI systems should encourage learning (reducing reformulation through informative responses) without creating dependency (over-responsive AI that accepts any query formulation reduces user incentive to learn effective communication).

\subsubsection{Future Directions}

Partial validation of bidirectional alignment components opens several research directions:

\begin{enumerate}
\item Complete coupling hypothesis test to test whether correlation between slope and diversity predicts helpfulness beyond individual metrics.
\item Semantic diversity reanalysis: Recompute responsiveness with SBERT embedding variance instead of Distinct-1.
\item Success stratification: Test whether negative reformulation slopes predict task success (learning interpretation) or failure (disengagement interpretation).
\item Causal intervention: A/B test with manipulated AI responsiveness levels.
\item Cross-dataset validation: Replicate findings on non-HH-RLHF datasets.
\item Bidirectional RLHF training: Implement training algorithms that optimize for coupling strength.
\end{enumerate}
